% LaTeX companion of hw03.typ. Both formats are editable.
\chapter{Homework 3: sequence limits and topology}

\begin{qlremark}{Authority label}

This is a visual transcription of \textbf{personal work} in \texttt{451-hw-3.pdf}. The raw assignment and consolidated solution file were used only for checking problem identities and notation; they are not the authority for any solution below.

\end{qlremark}

The submission begins with the definition: a sequence \(\left( a_{n} \right)\) of real numbers is \textbf{eventually constant} if there are \(c \in \mathbb{R}\) and \(N \in \mathbb{N}\) such that \(a_{n} = c\) for all \(n \geq N\).

\section{Problem 1 --- reciprocals and divergence to infinity}

Consider the bi-implication \(\lim a_{n} = \infty\Leftrightarrow\lim\frac{1}{a_{n}} = 0\).

\begin{qlsolution}{Forward direction}

Suppose \(\lim a_{n} = \infty\). Let \(\varepsilon > 0\) and consider \(M = \frac{1}{\varepsilon}\). Then for some \(N > M\), \(a_{n} > M\) whenever \(n \geq N\). Thus \(a_{n} > \frac{1}{\varepsilon}\) implies \(\left. |\frac{1}{a_{n}} \middle| < \varepsilon \right.\) whenever \(n \geq N\). So \(\lim\frac{1}{a_{n}} = 0\).

\end{qlsolution}

\begin{qlsolution}{Backward direction --- counterexample}

Consider \(a_{n} = - n\). Then \(\lim\frac{1}{a_{n}} = \lim\left( {- \frac{1}{n}} \right) = 0\), but \(\lim a_{n} = - \infty\).

\end{qlsolution}

\section{Problem 2 --- a bounded factor}

Let \(\left( a_{n} \right)\) and \(\left( b_{n} \right)\) be sequences of real numbers. Prove that if \(\lim a_{n} = 0\) and \(\left( b_{n} \right)\) is bounded, then \(\lim a_{n}b_{n} = 0\).

\begin{proof}

Since \(\left( b_{n} \right)\) is bounded, \(\left( {|b_{n}|} \right)\) is also bounded. Consider the constant sequence \(s_{n} = \sup\left( {|b_{n}|} \right)\). Then \(\left. \lim\left( s_{n} \right) = \sup \middle| b_{n}| \right.\). Since \(\lim a_{n} = 0\), \(\lim\left( {|a_{n}|} \right) = 0\), and hence \(\lim\left( |a_{n} \middle| \cdot \middle| s_{n}| \right) = 0\). As \(\left. s_{n} = \sup \middle| b_{n}| \right.\), \(\left. 0 \leq \middle| b_{n} \middle| \leq \middle| s_{n}| \right.\) for all \(n \in \mathbb{N}\), so \(\left. 0 \leq \middle| a_{n}b_{n} \middle| \leq \middle| a_{n}s_{n}| \right.\) for all \(n \in \mathbb{N}\). By the squeeze theorem, \(\left. \lim \middle| a_{n}b_{n} \middle| = 0 \right.\). Therefore \(\left. \lim a_{n}b_{n} = \lim \middle| a_{n}b_{n} \middle| = 0 \right.\).

\end{proof}

\section{Problem 3 --- three limits}

Determine the limits in the extended real line (including positive or negative infinity) of the following sequences and prove the results.

\subsection{\texorpdfstring{(a) \(\frac{2^{n}}{n!}\)}{(a) \textbackslash frac\{2\^{}\{n\}\}\{n!\}}}

\begin{proof}

The submitted answer is \(\lim_{n\rightarrow\infty}\frac{2^{n}}{n} \neq 0\). Let \(a_{n} = \frac{2^{n}}{n!}\). Then

\[\lim\limits_{n\rightarrow\infty}\frac{a_{n + 1}}{a_{n}} = \lim\limits_{n\rightarrow\infty}\frac{2^{n + 1}n!}{\left( {n + 1} \right)!2^{n}} = \lim\limits_{n\rightarrow\infty}\frac{2}{n + 1} = 0 < 1.\]

So \(\lim\frac{2^{n}}{n} \neq 0\).

\end{proof}

\subsection{\texorpdfstring{(b) \(\frac{n^{n}}{n!}\)}{(b) \textbackslash frac\{n\^{}\{n\}\}\{n!\}}}

\begin{proof}

The submitted answer is \(\lim_{n\rightarrow\infty}\frac{n^{n}}{n} \neq + \infty\). It writes \(\frac{n^{n}}{n} \neq \frac{n}{n - 1} \cdot \frac{n}{n - 2}\ldots\frac{n}{1} > n\). Let \(M > 0\) and choose an integer \(N \geq M\). For \(n \geq N\), \(\frac{n^{n}}{n!} > n \geq N > M\). Hence the limit is \(+ \infty\).

\end{proof}

\subsection{\texorpdfstring{(c) \(b_{1} = 2\), \(b_{n + 1} = \frac{b_{n}^{2} + 2}{2b_{n}}\)}{(c) b\_\{1\} = 2, b\_\{n + 1\} = \textbackslash frac\{b\_\{n\}\^{}\{2\} + 2\}\{2b\_\{n\}\}}}

\begin{proof}

Assume \(\lim b_{n} = L\). Then \(L = \lim b_{n + 1} = \lim\left( {\frac{b_{n}}{2} + \frac{1}{b_{n}}} \right) = \frac{L}{2} + \frac{1}{L}\), so \(\frac{L^{2}}{2} = 1\) and \(L = \sqrt{2}\). Since \(b_{n} > 0\) for all \(n \in \mathbb{N}\), the limit can only be \(\sqrt{2}\) if it exists.

Now prove \(\left( b_{n} \right)\) converges. For \(n \in \mathbb{N}\), \(b_{n + 1} = \frac{b_{n}}{2} + \frac{1}{b_{n}} \geq 2\sqrt{\frac{b_{n}}{2} \cdot \frac{1}{b_{n}}} = \sqrt{2}\). Since \(b_{1} = 2\), \(b_{n} \geq \sqrt{2}\) for all \(n \in \mathbb{N}\). Also \(\frac{b_{n + 1}}{b_{n}} = \frac{1}{2} + \frac{1}{b_{n}^{2}} \leq 1\). Hence \(\left( b_{n} \right)\) is decreasing and bounded below, so it converges. Therefore \(\lim b_{n} = \sqrt{2}\).

\end{proof}

\section{Problem 4 --- limits in a discrete set}

\begin{qlnote}{原稿红字}

hw 3 ①：discrete \(A \subseteq \mathbb{R}\) 中的任意 seq 要么 eventually constant，要么 \(\lim\left( a_{n} \right)\) 在 \(A\) 之外。

\end{qlnote}

Suppose \(A\) is a discrete subset of \(\mathbb{R}\), and \(\left( a_{n} \right)\) is a convergent sequence of numbers in \(A\). Prove that either \(\left( a_{n} \right)\) is eventually constant or \(\lim a_{n} \notin A\).

\begin{proof}

Write \(\lim a_{n} = L\). Assume \(\left( a_{n} \right)\) is not eventually constant and \(\lim a_{n} \in A\). Since \(A\) is discrete, there is some \(\varepsilon > 0\) such that \(\left( {L - \varepsilon,L + \varepsilon} \right) \cap A \smallsetminus \left\{ L \right\} = \varnothing\). Since \(\lim a_{n} = L\), there is \(N \in \mathbb{N}\) such that \(\left. |a_{n} - L \middle| < \varepsilon \right.\) for all \(n \geq N\). Since \(\left( a_{n} \right)\) is not eventually constant, there is \(n \geq N\) such that \(a_{n} \neq L\) and \(\left. |a_{n} - L \middle| < \varepsilon \right.\), i.e. \(a_{n} \in \left( {L - \varepsilon,L + \varepsilon} \right)\). Thus \(a_{n} \in \left( {L - \varepsilon,L + \varepsilon} \right) \cap A \smallsetminus \left\{ L \right\}\), a contradiction. Therefore \(\left( a_{n} \right)\) is either eventually constant or \(\lim a_{n} \notin A\).

\end{proof}

\section{Problem 5 --- sequences of rationals with bounded numerators}

For positive integer \(M\), let \(\mathbb{Q}_{M}\) be the set of rational numbers \(\frac{m}{n}\) with \(m,n \in \mathbb{Z}\) and \(\left. |m \middle| \leq M \right.\). Prove every sequence of distinct numbers in \(\mathbb{Q}_{M}\) converges.

\begin{proof}

Let \(\left( a_{n} \right)\) be an arbitrary sequence in \(\mathbb{Q}_{M}\) and let \(\varepsilon > 0\). Since for each \(q \in \mathbb{Z}\) there are only finitely many terms of \(\left( a_{n} \right)\) that have \(q\) as a denominator, consider \(N = \max\{ k:a_{k} = \frac{p}{q}\) for some \(p \leq M\) and \(q\) an integer with \(q \geq \frac{M}{\varepsilon}\}.\) Take arbitrary \(n \geq N + 1\). Then \(a_{n} = \frac{m}{q}\) where \(q > \frac{M}{\varepsilon}\). Thus \(a_{n} \leq \frac{M}{q} < \varepsilon\). So \(\lim a_{n} = 0\). This finishes the proof that every sequence of distinct numbers in \(\mathbb{Q}_{M}\) converges.

\end{proof}

\section{Problem 6 --- strict inequalities between sequences}

Let \(a_{n} < b_{n}\) for all \(n\).

\subsection{(a)}

\begin{proof}

Suppose \(\lim a_{n} = \infty\). Let \(M > 0\) and fix it. Then for some \(N \in \mathbb{N}\), \(a_{n} > M\) whenever \(n \geq N\). Since \(a_{n} < b_{n}\) for all \(n\), \(b_{n} > a_{n} > M\) for all \(n \geq N\). Therefore \(\lim b_{n} = \infty\).

\end{proof}

\subsection{(b)}

\begin{qlsolution}{Solution}

Consider \(a_{n} = \frac{1}{n^{2}}\) and \(b_{n} = \frac{2}{n^{2}}\) for all \(n \in \mathbb{N}\). Then \(a_{n} < b_{n}\) for all \(n \in \mathbb{N}\), but \(\lim a_{n} = \lim b_{n} = 0\).

\end{qlsolution}

\section{Problem 7 --- ratio limit greater than one}

\begin{qlnote}{原稿红字}

hw 3 ②：if positive seq \(\left( a_{n} \right)\) and \(\lim\frac{a_{n + 1}}{a_{n}} = L > 1\)，则 \(\lim\left( a_{n} \right) = \infty\)。

\end{qlnote}

Let \(\left( a_{n} \right)\) be a sequence of positive real numbers. Show that if \(\lim\frac{a_{n + 1}}{a_{n}} = L > 1\), then \(\lim a_{n} = \infty\).

\begin{proof}

Let \(\varepsilon = \frac{L - 1}{2}\). Since \(\lim\frac{a_{n + 1}}{a_{n}} = L\), there is some \(N_{1} \in \mathbb{N}\) such that \(\left. |\frac{a_{n + 1}}{a_{n}} - L \middle| < \varepsilon \right.\) for all \(n \geq N_{1}\), i.e. \(a_{n + 1} > \left( {\frac{L}{2} + \frac{1}{2}} \right)a_{n}\) for all \(n \geq N_{1}\). Let \(M > 0\). There is some \(N_{2} \geq N_{1}\) such that \(\left( {\frac{L}{2} + \frac{1}{2}} \right)^{N_{2}}a_{N_{2}} > M\), since \(\frac{L}{2} + \frac{1}{2} > 1\). Then for all \(n \geq N_{2}\), \(a_{n} \geq \left( {\frac{L}{2} + \frac{1}{2}} \right)^{n}a_{N_{2}} > M\). Therefore \(\lim a_{n} = \infty\).

\end{proof}

\section{Problem 8 --- lim sup and lim inf}

Find the lim sup and lim inf of the following sequences.

\begin{itemize}
\tightlist
\item
  (a) \(a_{n} = \left( {- 1} \right)^{n + 1} + \frac{\left( {- 1} \right)^{n}}{n}\): \(\operatorname{lim\, sup}\left( a_{n} \right) = 1\) and \(\operatorname{lim\, inf}\left( a_{n} \right) = - 1\).
\item
  (b) \(b_{n} = \sin\left( \frac{1}{n} \right)\): \(\operatorname{lim\, sup}\left( b_{n} \right) = \operatorname{lim\, inf}\left( b_{n} \right) = 0\).
\item
  (c) \(c:\mathbb{N}\rightarrow\mathbb{Q}\) any bijection: \(\operatorname{lim\, sup}\left( c_{n} \right) = + \infty\) and \(\operatorname{lim\, inf}\left( c_{n} \right) = - \infty\).
\item
  (d) \(d_{n} = \ln n + \cos n\): \(\operatorname{lim\, sup}\left( d_{n} \right) = \operatorname{lim\, inf}\left( d_{n} \right) = + \infty\).
\end{itemize}

\clearpage

\section{Problem 9 --- a recursive average}

Let \(a,b \in \mathbb{R}\) with \(a < b\). Let \(s_{1} = a\), \(s_{2} = b\), and \(s_{n + 2} = \frac{s_{n} + s_{n + 1}}{2}\). The submitted claim is \(\lim_{n\rightarrow\infty}s_{n} = \frac{2}{3}b + \frac{1}{3}a\).

\begin{proof}

Let \(d_{n} = s_{n + 1} - s_{n}\) for all \(n \in \mathbb{N}\). Then \(d_{1} = s_{2} - s_{1} = b - a\), and, for \(n \geq 2\),

\[d_{n} = \frac{s_{n - 1} + s_{n}}{2} - s_{n} = - \left( \frac{1}{2} \right)s_{n} - \left( \frac{1}{2} \right)s_{n - 1} = - \left( \frac{1}{2} \right)d_{n - 1}.\]

For all \(n \in \mathbb{N}\),

\[s_{n + 1} = \sum\limits_{i = 1}^{n{({s_{n + 1} - s_{n}})}} + s_{1} = s_{1} + \sum\limits_{i = 1}^{n}d_{n} = a + \frac{1 - \left( {- \frac{1}{2}} \right)^{n}}{1 - \left( {- \frac{1}{2}} \right)}d_{1} = a + \frac{2}{3}\left( {1 - \left( {- \frac{1}{2}} \right)^{n}} \right)\left( {b - a} \right).\]

Thus

\[\lim\limits_{n\rightarrow\infty}s_{n} = \lim\limits_{n\rightarrow\infty}s_{n + 1} = \lim\limits_{n\rightarrow\infty}\left( {a + \frac{2}{3}\left( {b - a} \right) - \frac{2}{3}\left( {- \frac{1}{2}} \right)^{n{({b - a})}}} \right) = \frac{2}{3}b + \frac{1}{3}a,\]

since \(\left. | - \frac{1}{2} \middle| < 1 \right.\).

\end{proof}

\section{Problem 10 --- a divergent sequence with one possible subsequential limit}

Consider \(a_{n} = n^{{({- 1})}^{n}}\), i.e. \(\left( a_{n} \right) = \left( {1,2,\frac{1}{3},4,\frac{1}{5},6,\ldots} \right)\). The work claims \(\left( a_{n} \right)\) diverges, but every convergent subsequence converges to \(L = 0\).

\begin{proof}

The odd-indexed terms \((a_{n_{k}}:k\) is odd\() = \left( {1,\frac{1}{3},\frac{1}{5},\ldots} \right)\rightarrow 0\). Let \(\left( a_{n_{k}} \right)\) be a convergent subsequence of \(\left( a_{n} \right)\); then \(k\mapsto n_{k}\) is strictly increasing. Suppose there are infinitely many \(k \in \mathbb{N}\) such that \(n_{k}\) is even. We show \(\left( a_{n_{k}} \right)\) diverges. Let \(L \in \mathbb{R}\), take \(M = 1\), and fix \(N \in \mathbb{N}\). If there is no \(n_{k} > N\) with \(a_{n_{k}} > L + 1\), then there are only finitely many even \(n_{k}\), a contradiction. Thus there must be \(n_{k} > N\) with \(\left. |a_{n_{k}} - L \middle| > M \right.\), so \(\left( a_{n_{k}} \right)\) diverges. Hence only finitely many \(n_{k}\) are even. Cutting the tail makes all remaining \(n_{k}\) odd, and so \(\left( a_{n_{k}} \right)\) converges to \(0\).

\end{proof}

\section{Problem 11 --- lim sup of a sum}

Let \(\left( a_{n} \right)\) and \(\left( b_{n} \right)\) be bounded sequences of positive real numbers.

\subsection{(a)}

\begin{proof}

Write \(l_{n} = \sup\left\{ {a_{k} + b_{k}:k \geq n} \right\}\), \(u_{n} = \sup\left\{ {a_{k}:k \geq n} \right\}\), and \(v_{n} = \sup\left\{ {b_{k}:k \geq n} \right\}\). Let \(\varepsilon > 0\). Then for every \(k \geq n\), \(a_{k} < u_{n} + \frac{\varepsilon}{2}\) and \(b_{k} < v_{n} + \frac{\varepsilon}{2}\), so \(a_{k} + b_{k} < u_{n} + v_{n} + \varepsilon\). Hence \(l_{n} \leq u_{n} + v_{n}\). Since \(n\) is arbitrary, \(\lim l_{n} \leq \lim u_{n} + \lim v_{n}\), i.e. \(\operatorname{lim\, sup}\left( {a_{n} + b_{n}} \right) \leq \operatorname{lim\, sup}\left( a_{n} \right) + \operatorname{lim\, sup}\left( b_{n} \right)\).

\end{proof}

\subsection{(b)}

\begin{qlsolution}{Counterexample}

\(a_{n} = 1 + \left( {- 1} \right)^{n}\), so \(\operatorname{lim\, sup}\left( a_{n} \right) = 2\). Let \(b_{n} = 1 + \left( {- 1} \right)^{n + 1}\), so \(\operatorname{lim\, sup}\left( b_{n} \right) = 2\). But \(\operatorname{lim\, sup}\left( {a_{n} + b_{n}} \right) = 1 + 1 = 2 < \operatorname{lim\, sup}\left( a_{n} \right) + \operatorname{lim\, sup}\left( b_{n} \right)\).

\end{qlsolution}

\subsection{(c)}

\begin{qlsolution}{Solution}

Write \(\lim a_{n} = L\). Then \(\operatorname{lim\, sup}\left( a_{n} \right) = \lim a_{n} = L\) since \(\left( a_{n} \right)\) converges. Thus \(\operatorname{lim\, sup}\left( a_{n} \right) + \operatorname{lim\, sup}\left( b_{n} \right) = L + \lim v_{n} = \lim\left( {L + v_{n}} \right) = \operatorname{lim\, sup}\left( {a_{n} + b_{n}} \right)\).

\end{qlsolution}

\section{Problem 12 --- a sequence with every real subsequential limit}

\begin{qlnote}{原稿红字}

hw 3 ③：\(\mathbb{R}\) 中存在一个 seq，使得 ``all sub seq lim of \(\left( a_{n} \right)\)'' \(= \mathbb{R}\)。

\end{qlnote}

\begin{proof}

Since \(\mathbb{N} \approx \mathbb{Q}\), there exists a surjective function \(S:\mathbb{N}\rightarrow\mathbb{Q}\). Note that \(\left( S_{n} \right)\) is a sequence. Let \(r \in \mathbb{R}\) be arbitrary. There exists a sequence in \(\mathbb{Q}\), \(\left( q_{n} \right)\rightarrow r\). Since \(S:\mathbb{N}\rightarrow\mathbb{Q}\) is surjective, consider the subsequence \(\left( S_{n_{k}} \right)\) of \(\left( S_{n} \right)\) defined by \(S_{n_{k}} = q_{n}\) for some \(n \in \mathbb{N}\), for all \(k \in \mathbb{N}\). Take a monotonic subsequence of \(\left( S_{n_{k}} \right)\) as \(\left( S_{m} \right)\). It is a subsequence of \(\left( S_{n_{k}} \right)\), so it is also a subsequence of \(\left( S_{n} \right)\). Let \(\varepsilon > 0\). There is some \(N \in \mathbb{N}\) such that \(\left. |q_{n} - r \middle| < \varepsilon \right.\) whenever \(n \geq N\). Since there is some term \(S_{m}\) with \(S_{m} = q_{N}\) and \(\left( S_{m} \right)\) is monotonic, \(\left. |S_{m} - r \middle| < \varepsilon \right.\) whenever \(m \geq M\). Therefore \(\left( S_{m} \right)\rightarrow r\).

\end{proof}

\section{Problem 13 --- open and closed sets}

The submitted classifications are:

\begin{itemize}
\tightlist
\item
  (a) \(\left\{ {\frac{1}{n}:n \in \mathbb{N}} \right\}\): neither.
\item
  (b) \(\left\{ {\frac{1}{n}:n \in \mathbb{N}} \right\} \cup \left\{ 0 \right\}\): closed and not open.
\item
  (c) \(\cup_{n \geq 1}\left\lbrack {\frac{1}{n},3 - \frac{1}{n}} \right\rbrack\): open and not closed.
\item
  (d) \(\mathbb{Z}\): closed and not open.
\item
  (e) \(\mathbb{Q}\): neither.
\item
  (f) \(\cap_{n \geq 1}\left( {- \frac{1}{n},\frac{1}{n}} \right)\): closed and not open.
\end{itemize}

\section{Problem 14 --- closed discrete set with no uniform separation}

\begin{qlsolution}{Counterexample}

Consider \(S_{n} = \sum_{k = 1}^{n}\frac{1}{k}\), a partial sum of the harmonic series, and \(A = \left\{ {S_{n}:n \in \mathbb{N}} \right\}\). There is no subsequential limit in \(S_{n}\), so \(A\) has no limit point; hence \(A = A'\) and \(A\) is closed. For each \(S_{n}\) consider \(\varepsilon = \frac{1}{n + 1}\); then \(V_{\varepsilon{(S_{n})}} \cap A \smallsetminus \left\{ S_{n} \right\} = \varnothing\), so \(A\) is discrete. But there is no \(\varepsilon > 0\) such that \(\left. |a - b \middle| \geq \varepsilon \right.\) for every pair \(a,b \in A\), since for any \(\varepsilon > 0\), \(S_{k + 1} - S_{k} < \frac{1}{\varepsilon} = \varepsilon\) (as recorded in the submission).

\end{qlsolution}

\section{Problem 15 --- an external limit point}

\begin{qlnote}{原稿红字}

hw 3 ④：bounded + infinite + discrete \(A \subseteq \mathbb{R}\) 一定存在不在 \(A\) 中的 subseq lim。

\end{qlnote}

Suppose \(A \subseteq \mathbb{R}\) is infinite, bounded, and discrete. Prove that there is a convergent sequence in \(A\) whose limit is not in \(A\).

\begin{proof}

Take an arbitrary sequence \(\left( a_{n} \right)\) in \(A\) such that \(\forall m,n \in \mathbb{N}\), \(a_{m} \neq a_{n}\). By the Bolzano--Weierstrass theorem, there is a convergent subsequence \(\left( a_{n_{k}} \right)\); write \(\lim a_{n_{k}} = L\). Claim: \(L \notin A\). Suppose \(L \in A\). Since \(L\) is the limit of a sequence in \(A\), it is a limit point of \(A\), so for every \(\varepsilon > 0\) there exists \(x \in A \smallsetminus \left\{ L \right\}\) with \(\left. 0 < \middle| x - L \middle| < \varepsilon \right.\), i.e. \(x \in V_{\varepsilon{(L)}} \cap A \smallsetminus \left\{ L \right\}\). Since \(A\) is discrete and \(L \in A\), there exists \(\varepsilon > 0\) such that \(V_{\varepsilon{(L)}} \cap A = \left\{ L \right\}\). These two statements contradict. So \(L \notin A\).

\end{proof}
